\documentclass[12pt]{article}
\usepackage[a4paper, margin=1in]{geometry}
\usepackage{enumitem}

\begin{document}

\section*{CSE 400 L4 Scribe}

Generate a lecture scribe intended to serve as exam-oriented reference material for CSE 400.

Use only the provided context (lecture slides and relevant textbook material). Do not introduce any material, definitions, examples, explanations, assumptions, intuition, or notation that are not explicitly present in the context.

While writing the scribe, reproduce the lecture content faithfully and in the same scope and structure as taught. Where the context contains definitions, assumptions, propositions, theorems, proofs, derivations, or worked examples, reconstruct them step by step exactly as they appear, making all logical steps explicit so that the reasoning can be followed during exam revision. Do not add reasoning where none is present in the context.

Include proofs, derivations, and worked examples only if they are present in the provided context, and reproduce them without modification, simplification, or re-interpretation.

Organize the scribe clearly using appropriate sectioning and headings so that it can be reliably used for closed-notes exam preparation.

Prioritize correctness, faithfulness to the lecture and textbook materials, and explicit logical flow. Do not invent or infer content beyond the provided context.

\bigskip

\textbf{this is the context of L4:}

For Lecture--4 of CSE 400 scheduled on 15th January, the class will proceed according to the following agenda:

\begin{enumerate}
\item Attendance  
Attendance will be taken at the beginning of the class. Please ensure punctuality.

\item Motivation through Engineering Applications
\begin{itemize}
\item Speech Recognition
\item Radar Systems
\item Communication Networks
\end{itemize}

\item Introduction to Probability Theory
\begin{itemize}
\item Random experiments, sample space, and events
\item Axioms of probability
\item Corollaries and propositions derived from the axioms
\item Assigning probability: classical and relative frequency approaches
\end{itemize}

\item Joint Probability
\begin{itemize}
\item Motivation, notation, and key ideas
\item Examples: Card deck example, Costume party example
\end{itemize}

\item Conditional Probability
\begin{itemize}
\item Motivation, notation, and core concepts
\item Examples: Cards without replacement, Game of Poker, The missing key
\end{itemize}

\item Graded Class Participation  
Class participation based on the content covered during the lecture will be assessed.
\end{enumerate}

\bigskip

Below is a faithful, exam-oriented lecture scribe for CSE 400 -- Lecture 4 (15 January), written strictly using only the provided lecture context.  
No definitions, explanations, examples, notation, assumptions, or derivations beyond what is explicitly listed in the context have been added.

\section*{CSE 400 -- Lecture 4}

\textbf{Date:} 15 January  
\textbf{Lecture Scribe (Exam-Oriented Reference)}

\section*{1. Attendance}
Attendance was taken at the beginning of the class.  
Students were required to ensure punctuality.

\section*{2. Motivation through Engineering Applications}
The lecture began with motivation for the topic using the following engineering applications:

\begin{itemize}
\item Speech Recognition
\item Radar Systems
\item Communication Networks
\end{itemize}

\section*{3. Introduction to Probability Theory}

\subsection*{3.1 Random Experiments, Sample Space, and Events}
\begin{itemize}
\item Random experiments
\item Sample space
\item Events
\end{itemize}

\subsection*{3.2 Axioms of Probability}
\begin{itemize}
\item Axioms of probability
\end{itemize}

\subsection*{3.3 Corollaries and Propositions Derived from the Axioms}
\begin{itemize}
\item Corollaries derived from the axioms of probability
\item Propositions derived from the axioms of probability
\end{itemize}

\subsection*{3.4 Assigning Probability}
\begin{itemize}
\item Classical approach to assigning probability
\item Relative frequency approach to assigning probability
\end{itemize}

\section*{4. Joint Probability}

\subsection*{4.1 Motivation, Notation, and Key Ideas}
\begin{itemize}
\item Motivation for joint probability
\item Notation used for joint probability
\item Key ideas related to joint probability
\end{itemize}

\subsection*{4.2 Examples}
\begin{itemize}
\item Card deck example
\item Costume party example
\end{itemize}

\section*{5. Conditional Probability}

\subsection*{5.1 Motivation, Notation, and Core Concepts}
\begin{itemize}
\item Motivation for conditional probability
\item Notation used in conditional probability
\item Core concepts of conditional probability
\end{itemize}

\subsection*{5.2 Examples}
\begin{itemize}
\item Cards without replacement
\item Game of Poker
\item The missing key
\end{itemize}

\section*{6. Graded Class Participation}
Class participation was assessed.  
Assessment was based on the content covered during the lecture.

\bigskip

\textbf{End of Lecture 4 Scribe}

\bigskip

If you want, once you share the actual slide content or textbook excerpts, I can convert this into a fully detailed, proof-by-proof, exam-ready scribe without violating your constraints.

\end{document}
